(1) 两个无自旋玻色子局限在两个站点，哈密顿量 \(H=-J(a_1^\dagger a_2+a_2^\dagger a_1)+\frac U2[n_1(n_1-1)+n_2(n_2-1)]\)，\(J>0,\ U\ge0\)，总数 \(N=2\)。在正交Fock基 \(|20\rangle,|11\rangle,|02\rangle\) 中写出完整矩阵，并用交换对称/反对称组合严格对角化，给出三条本征能量。
(2) 初态为 \(|11\rangle\)。求两粒子占同一站点的总概率 \(P_{\rm double}(t)\)，及每一站点平均占据数，给出 \(U=0\) 与 \(U\gg J\) 极限和振荡时间尺度。
(3) 在 \(U/J\gg1\) 的排斥区对最低能量展开至 \(J^2/U\)，解释虚双占据所产生的能量下降，并检查该能量修正与二阶微扰理论中两条独立中间跃迁路径相符。
(4) 求基态在站点1/站点2二分下的约化密度矩阵本征概率，写出站点纠缠的von Neumann熵。分别求 \(U=0\) 与 \(U/J\to\infty\) 的精确值，并解释这里的站点纠缠不等同于为两个可区别粒子定义的纠缠。
